%% ma: L11-B06-0018
%% lop: 11
%% chuong: 2
%% bai: 6
%% dang: 11.6.12
%% loai: TL
%% muc: TH
%% dap_an: "-"
%% nguon: "de-on-tap-giua-ki-1-toan-11-nam-2026-2027-truong-thpt-ngo-quyen-da-nang.pdf (M09, Đề ôn tập giữa kì 1 Toán 11 THPT Ngô Quyền, Đà Nẵng), Mã đề 0111, Phần IV Câu 2"
%% kien_thuc: "Bài 6 (cấp số cộng)"
%% trang_thai: cho-duyet
%% ly_do: "Dạng mới đề xuất 11.6.12"
\begin{ex}
Một vận động viên điền kinh chuẩn bị tham gia giải chạy marathon. Để rèn luyện thể lực, vận động viên lập kế hoạch luyện tập trong $30$ ngày liên tục. Ngày thứ nhất, vận động viên chạy cự li $4$ km. Kể từ ngày thứ hai trở đi, mỗi ngày vận động viên đều chạy nhiều hơn ngày liền trước đó $500$ m.

a) Tính cự li mà vận động viên đó chạy trong ngày thứ $15$ của kế hoạch luyện tập.

b) Tính tổng quãng đường mà vận động viên đã chạy sau khi hoàn thành toàn bộ kế hoạch $30$ ngày luyện tập.
\loigiai{Cự li các ngày lập thành cấp số cộng $u_1=4$, $d=0{,}5$ (km).

a) $u_{15}=4+14\cdot0{,}5=11$ (km).

b) $S_{30}=\dfrac{30}{2}(2\cdot4+29\cdot0{,}5)=15\cdot22{,}5=337{,}5$ (km).}
\end{ex}
